%%%%%%%%%%%%%%%%%%%%%%%%%%%%%%%%%%%%%%%%%%%%%%%%%%%%
% Template Artigo
%
%Template adotado a partir do modelo SBC, foram incluídos alguns pacotes extras e alguns ítens não contemplando no modelo da SBC.
%
% Marque esta página e compartilhe a URL com seus co-autores. Eles podem edite ao mesmo tempo!
%
% Você pode fazer upload de figuras, bibliografias, classes personalizadas e estilos usando o menu de arquivos.
%
%Autor: Rodrigo Rosa da Silva
%%%%%%%%%%%%%%%%%%%%%%%%%%%%%%%%%%%%%%%%%%%%%%%%%%%
\documentclass[12pt]{article}

\usepackage{sbc-template}
\usepackage{graphicx,url}
\usepackage{xcolor}
%\usepackage[brazil]{babel}   
\usepackage[utf8]{inputenc}
\usepackage[portuguese]{babel}
\graphicspath{ {./} }
\usepackage{quoting}
\usepackage{hyphenat}
\usepackage{hyperref}

\sloppy
\title{Título do Artigo e Subtítulo se for o Caso}

\author{Fabiane N. Prates\inst{1}, Carlos E. P. Severo\inst{1}, Rodrigo R. Silva\inst{2} }

\address{Instituto Federal de Educação, Ciência e Tecnologia Sul-rio-grandense (IFSul)\\
  Av. Leonel de Moura Brizola, 2501 -- 96418-400 -- Bagé -- RS -- Brasil
\nextinstitute
  Prefeitura Municipal de Bagé\\
  Av. Gen. Osório, 998 -- 96400-100 -- Bagé - -- RS -- Brasil 
    \email{\{fabiane,calos\}@ifsul.edu.br,rodrigo@gmail.com}
}

\begin{document} 

\maketitle

\begin{abstract}
  Considering (Justification / problem), objective (Objectives). For that, proceeded to (Methodology). Thus, it is observed that (Results), which allows to conclude that (Conclusion).
  
  \noindent
  \textbf{Keywords}: Technology, Android, Data Mining.
\end{abstract}

\begin{resumo} 
  Considerando (Justificativa/problema), objetiva-se(Objetivos). Para tanto procedeu-se à (Metodologia). Desse modo, observa-se que (Resultados), o que permite concluir que (Conclusão).
  
  \noindent
  \textbf{Palavras-chave}: Tecnologia, Android, Mineração de dados.
\end{resumo}


\section{Introdução}

O formato do papel deve ser A4 com coluna única, 3,5 cm para superior margem, 2,5 cm para a margem inferior e 3,0 cm para as margens laterais, sem cabeçalhos ou rodapés. A fonte principal deve ser Times, tamanho nominal 12 pontos, com 6 pontos de espaço antes de cada parágrafo. Abaixo uma sugestão de divisão do artigo em sessões.

\section{Trabalhos Correlatos} \label{sec:firstpage}
No âmbito da disponibilização de informações geoespaciais, foram identificados trabalhos relacionados ao problema investigado nesta pesquisa...

Sugestão de 3 a 4 trabalhos com referências, com objetivo, metodologia, resultado e conclusão sobre cada um. Ao final, comentar sobre os trabalhos e a diferença entre a proposta...


\section{Metodologia}
 Pode-se usar subseções no artigo para melhor delinear as etapas.

\subsection{Imagens}
Nesta sessão apresenta-se algumas maneiras de incluir imagens. Procure utilizar sempre o formato PDF para não perder a qualidade vetorial.

Abaixo o código de uso para uma imagem, no caso a Figura \ref{fig:Tipos de Solos}, use o label para referenciar as imagens e gerar a numeração automática. Na fonte também pode-se usar referências do .bib conforme Figura \ref{fig:modeloCascata}. As imagens estão dentro de um diretório chamado \textit{imagens} para fins de organização...
\begin{figure}[ht]
    \centering
    \includegraphics[width=.75\textwidth]{imagens/Rplot.pdf}\\
    Fonte: Próprio Autor
    \caption{Tipos de Solos}
    \label{fig:Tipos de Solos}
\end{figure}

\begin{figure}[ht]
    \centering
    \includegraphics[width=0.75\textwidth]{imagens/Cascata.pdf}\\
    Fonte:\cite{stankiewicz2017modelo}
    \caption{O modelo em cascata}
    \label{fig:modeloCascata}
\end{figure}


\newpage%Incluido para quebrar página, no artigo retire...

\subsection{Tabelas}
Existem muitas formas de se criar tabelas no \LaTeX, sugere-se um editor online para propor agilidade.
Sugestão \href{https://www.tablesgenerator.com/}{EDITOR ONLINE}. É possível criar uma tabela totalmente personalizada com linhas, colunas, cores etc. A Tabela \ref{tab:tabelaModelo} é apenas um modelo ilustrativo...

\begin{table}[]
    \centering
    \caption{Tabela de Números}\vspace{0.4cm}
    \begin{tabular}{|c|c|c|}
    \hline
    \textbf{Número} & \textbf{Dobro} & \textbf{Triplo} \\ \hline
    20  & 40  & 60 \\ \hline
    19  & 38  & 57 \\ \hline
    18  & 36  & 54 \\ \hline
    ... & ... & ...\\ \hline
    1   & 1   & 1  \\ \hline
    \end{tabular}
    \vspace{0.4cm}
    \\Fonte: Próprio Autor
    \label{tab:tabelaModelo}
\end{table}
\newpage%Incluido para quebrar página, no artigo retire...
\subsection{Citações}
As referências bibliográficas devem ser inequívocas e uniformes. Recomenda-se fornecer referências aos nomes dos autores entre colchetes, por exemplo, \cite{knuth:84}, \cite{boulic:91}. As referências devem ser listadas usando tamanho de fonte 12 pontos, com 6 pontos de espaço antes de cada referência. A primeira linha de cada referência não deve ser indentada, enquanto a seguinte deve ser indentada em 0,5 cm.

Se quiser que uma referência apareça listada nas referências mas não no texto, use nocite\{\} \nocite{smith:99}. As citações devem estar formatadas no arquivo sbc-template.bib. Sugestação, em pesquisa no google acadêmico e na maioria dos repositórios de artigos e trabalhos como {\it Web of Science} e {\it SCIELO} já se encontra a citação formatada para o padrão .bib, bastando copiar em colocar no arquivo sbc-template.bib. Este arquivo .bib já apresenta alguns formatos de referências para serem usados como base, por exemplo, leis, sites, artigos, livros etc...
\section{Resultados e Discussões}



\section{Conclusão}



\bibliographystyle{sbc}
\bibliography{sbc-template}

\end{document}
